\documentclass[10pt]{article}
\usepackage{titlesec}
\usepackage{geometry}
\geometry{verbose,tmargin=.9in,bmargin=.9in,lmargin=1.0in,rmargin=1.0in}
\usepackage{amsmath,amsfonts,amsthm,amssymb}
\usepackage{url}
\usepackage{color}
\usepackage[usenames,dvipsnames,svgnames,table]{xcolor}
\usepackage[colorlinks=true, linkcolor=red, urlcolor=blue, citecolor=gray]{hyperref}
\usepackage{float}
\usepackage{caption}
\usepackage{subcaption}
\usepackage{graphicx}
\usepackage{wrapfig}
\usepackage{booktabs}
\usepackage{longtable}
\usepackage{enumitem}
\usepackage{multicol}
\usepackage{etoolbox}
\usepackage{hyperref}
\newcommand{\bs}[1]{\boldsymbol{#1}}
\newcommand{\bv}[1]{\mathbf{#1}}
\newcommand{\R}{\mathbb{R}}
\newcommand{\E}{\mathbb{E}}
\DeclareMathOperator*{\Var}{var}
\DeclareMathOperator*{\argmin}{arg\,min}
\DeclareMathOperator*{\argmax}{arg\,max}

\definecolor{nyuDarkPurple}{HTML}{330662}
\definecolor{nyuOfficialPurple}{HTML}{57068c}

\newcommand{\spara}[1]{\vspace{.5em}\noindent {\large\sffamily\textcolor{nyuOfficialPurple}{#1}}}
\titleformat{\section}[hang]{\Large\sffamily\color{nyuDarkPurple}}{\thesection}{1em}{}
\titleformat{\subsection}[hang]{\large\sffamily\color{nyuDarkPurple}}{\thesection}{1em}{}
\titleformat{\subsubsection}[hang]{\normalsize\sffamily\color{gray}}{\thesection}{1em}{}

\usepackage{fancyhdr}
\pagestyle{fancy}
\lhead{\includegraphics[width=4cm]{tandon_long_color.eps}}
\rhead{\thepage}
\pagenumbering{gobble}

\setcounter{secnumdepth}{0}

\begin{document}
	
\begin{center}
	\normalsize
	New York University Tandon School of Engineering
	
	Computer Science and Engineering
	\medskip
	
	\large
	CS-GY 6763: Midterm Practice. 
	\medskip
\end{center} 

\section{Logistics}
\begin{itemize}
	\item Exam will be held in class on \textbf{Friday, 3/14} starting at 2:00pm sharp. Please arrive on time!
	\item You will have \textbf{1 hour, 15 minutes} to answer a variety of short answer and longer form questions.
	\item You can bring a one page sheet of paper (two-sided if you want) with notes, theorems, etc. written down for reference. 
	\item Scrap paper will be provided and can be used to write solutions if extra space is needed. 
	\item I will be in the room to answer any questions.
\end{itemize}

\section{Concepts to Know}

\subsection{Random variables and concentration.}
\begin{itemize}
	\item Linearity of expectation and variance.
	\item Indicator random variables and how to use them.
	\item Markov's inequality, Chebyshev's inequality (ideally should know from memory so you can apply quickly).
	\item Union bound  (should know from memory).
	\item Chernoff and Bernstein bounds (don't need to memorize the exact bounds, but can apply if given).
	\item General idea of law of large numbers and central limit theorem. 
	\item The probability that a normal random variables $\mathcal{N}(0,\sigma^2)$ falls further than $k\sigma$ away from its expectation is $\leq O(e^{-k^2/2})$. 
\end{itemize}

\subsection{Hashing, Dimensionality Reduction, High Dimensional Vectors}
\begin{itemize}
	\item Random hash functions.
	\item Random hashing for frequency estimation. 
	\item Random hashing for distinct elements estimation.
	\item MinHash for Jaccard similarity estimation. 
	\item Locality sensitive hash functions. 
	\item MinHash and SimHash for Jaccard Similarity and Cosine Similarity.
	\item Adjusting false positive rate and false negative rate in an LSH scheme. 
	\item Statement of Johnson-Lindenstrauss lemma (know from memory).
	\item Statement of \emph{distributional} JL lemma and how it can be used to prove JL.
\end{itemize}

\subsection{High dimensional geometry}
\begin{itemize}
	\item How to draw a random unit vector from the sphere in $d$ dimensions (draw $\bv{x}$ with all entries i.i.d. $\mathcal{N}(0,1)$ and normalize it). 
	\item How does $\|x - y\|_2^2$ relate to $\langle x, y\rangle$ if $x$ and $y$ are unit vectors?
%	\item Why is the inner product $\langle x,  y\rangle$ of two uniformly random unit vectors on the sphere identically distributed to to $\langle e_i,  y\rangle$ for any standard basis vector $e_i$? Or for that matter to $\langle z,  y\rangle$ for any fixed vector $z$?
	\item How many mutually orthogonal unit vectors are there in $d$ dimensions?
	\item There are $2^{\theta(\epsilon^2d)}$ nearly orthogonal unit vectors in $d$ dimensions (with $\langle x,  y\rangle \leq \epsilon$). Know roughly how prove this fact using the \emph{probabilistic method}, which required a an exponential \emph{concentration inequality} + \emph{union bound}.
	\item Know how to prove that all but an $2^{\theta(-\epsilon d)}$ fraction of a balls volume in $d$ dimensions lies in a spherical shell of width $\epsilon$ near its surface.
	\item The surface area/volume ratio \emph{increases} in high dimensions. 
	\item The cube volume/ball volume ratio \emph{increases} in high dimensions.
	
\end{itemize}

% \subsection{Convex optimization}
% \begin{itemize}
% 	\item Definition(s) of convex function.
% 	\item Definition of convex set.
% 	\item Gradient descent basic update rule. 
% 	\item Definitions of $G$-Lipschitz, $\beta$-smooth, $\alpha$-strongly convex. Know the first order definition for high-dimensional functions. The second order definition you only need to know for low-dimensional functions. I.e. a twice differentiable function $f:\R\rightarrow \R$ is $\beta$-smooth, $\alpha$-strongly convex if for all $x$, $\alpha \leq f''(x) \leq \beta$. I won't test on Hessians.
% 	\item Definition of condition number.
% 	\item How much time does it take to multiply an $n\times d$ matrix by a $d \times m$ matrix?
% 	\item Be able to compute gradients of basic functions from $\R^d \rightarrow \R$. 
% 	\item Definition(s) of convex function.
% \end{itemize}

	
	\section{Practice Problems}
		
	\begin{enumerate}
		\item Show that for any  random variable ${X}$, $\E[{X}^2] \ge \E[{X}]^2$.
		\item Show that for independent ${X}$ and ${Y}$ with $\E[{X}] = \E[{Y}] = 0$, $\Var[{X} \cdot {Y}] = \Var[X] \cdot \Var[Y]$. 
		\item Given a random variable $X$, can we conclude that $\E[1/X] = 1/E[X]$? If so, prove this. If not, give an example where the equality does not hold.
		\item Indicate whether each of the following statements is \textbf{always} true, \textbf{sometimes} true, or \textbf{never} true. Provide a short justification for your choice. 
		\begin{enumerate}[label=(\alph*)]
			\item $\Pr[{X} = s \text{ and } {Y} = t] > \Pr[{X} = s]$. \hspace{1em}ALWAYS\hspace{1em} SOMETIMES\hspace{1em} NEVER
			\item $\Pr[{X} = s \text{ or }  {Y} = t] \le \Pr[{X} = s] + \Pr[{Y} = t]$. \hspace{1em}ALWAYS \hspace{1em} SOMETIMES\hspace{1em} NEVER
			\item $\Pr[{X} = s \text{ and } {Y} = t] = \Pr[{X} = s] \cdot \Pr[{Y} = t]$. \hspace{1em}ALWAYS\hspace{1em} SOMETIMES \hspace{1em} NEVER
		\end{enumerate}
		\item Let $\bs{\Pi}$ be a random Johnson-Lindenstrauss matrix (e.g. scaled random Gaussians) with $O(\log(1/\delta)/\epsilon^2)$ rows. With probability $(1-\delta)$,
		\begin{align*}
			\min_{\bv{x}} \|\bs{\Pi}\bv{A}\bv{x} - \bs{\Pi}\bv{b}\|_2^2 \leq (1+\epsilon) \min_{\bv{x}} \|\bv{A}\bv{x} - \bv{b}\|_2^2. 
		\end{align*}
		Is the statement true: \hspace{1em}ALWAYS\hspace{1em} SOMETIMES \hspace{1em} NEVER

		\item Assume there are $1000$ registered users on your site $u_1, \ldots, u_{1000}$, and in a given day, each user
		visits the site with some probability $p_i$. The event that any user visits the site is independent of what the other users do. Assume that $\sum_{i=1}^{1000} p_i = 500$. 
		\begin{enumerate}
			\item  Let $X$ be the number of users that visit the site on the given day. What is $E[X]$?
			\item Apply a Chernoff bound to show that $Pr[X \geq 600] \leq .01$.
			\item Apply Markov’s inequality and Chebyshev’s inequality to bound the same probability.
			How do they compare?
		\end{enumerate}
		\item Give an example of a random variable and a deviation t where Markov’s inequality gives a	tighter upper bound than Chebyshev’s inequality.

%		\item SimHash can be applied to binary vectors ${y}, {q}\in \{0,1\}^d$. How does doing so compare to MinHash? How is the Jaccard similarity $J(y,q)$ related to the cosine similarity $\cos(\theta(y,q))$?
		\item Suppose there is some unknown vector $\bs{\mu} \in \R^d$. We receive noise perturbed random samples of the form $\bv{Y}_1 = \bs{\mu} + \bv{X}_1, \ldots, \bv{Y}_k = \bs{\mu} + \bv{X}_k$ where each $\bv{X}_i$ is a random vector with each of its entries distributed as an independent random normal $\mathcal{N}(0,1)$. From our samples $\bv{Y}_1, \ldots, \bv{Y}_k$ we hope to estimate $\bs{\mu}$ by $\bs{\tilde{\mu}} = \frac{1}{k}\sum_{i=1}^k \bv{Y}_i$. 
		\begin{enumerate}
			\item Prove that with $k = O(\log n /\epsilon^2)$ samples, $\max_{i=1,\ldots, d}|\bs{\mu}_i -\bs{\tilde{\mu}_i} | \leq \epsilon$ with probability $9/10$.
			\item Prove that with $k = O(n\log n /\epsilon)$ samples, $\|\bs{\mu} -\bs{\tilde{\mu}}\|_2 \leq \epsilon$ with probability $9/10$ \textbf{Bound}
		\end{enumerate}
%			Is the following also true with high probability?
%	\begin{align*}
%		(1-\epsilon) \min_{\bv{x}} \|\bv{A}\bv{x} - \bv{b}\|_2^2\leq \min_{\bv{x}} \|\bs{\Pi}\bv{A}\bv{x} - \bs{\Pi}\bv{b}\|_2^2
%	\end{align*}

% \item In Homework 1, we considered a random walk on a $d$-dimensional grid. At each step, the walk chooses one of the $d$-dimensions uniformly at random, and takes a step in that direction - up with probability $1/2$ and down with probability $1/2$. Assume that the walk starts at the origin, takes $n$ steps, and ends at position $x$, where $x$ is a $d$ dimensional vector with integer values. We proved that $\E[\|x\|_2^2] = n$. 
% \\\\
% Since the result does not depend on $d$, I said this fact was surprising: if a tourist starts at Washington Square Park and wonders around Manhattan (i.e., takes a two-dimensions random walk), in expectation they don't get any more ``lost'' than if they restrict their wondering to just up and down 5th Avenue (i.e., a one dimensional random walk.)
% \\\\
% However, it was pointed out by a student that this is not quite correct. To return back to the park, the tourist must walk along streets and avenues, so we should really use $\ell_1$ norm $\|x\|_1 = \sum_{i=1}^n |x_i|$ to determine how lost the tourist is. 

\item For two length $d$ binary vectors $\bv{q},\bv{y} \in \{0,1\}^d$, consider the hamming similarity: 
\begin{align*}
	s(\bv{q},\bv{y}) = 1 - \frac{\|\bv{q} - \bv{y}\|_0}{d}.
\end{align*} 
Construct a function $h$ as follows:
	define the random function $c: \{0,1\}^d \rightarrow \{0,1\}$ as $c(\bv{x}) = \bv{x}[j]$, where $j$ is a uniform random integer in $\{1, \ldots, d\}$. Then, let $g$ be a uniform random hash function from $\{0,1\} \rightarrow \{1, \ldots, m\}$. Finally, let:
	\begin{align*}
		h(\bv{x}) = g(c(\bv{x})).
	\end{align*}
	
	 Prove that $h$ is a locality sensitive hash function for {hamming similarity}. 
	\end{enumerate}
	
	
	
	% \subsection{Convex optimization}
	% From \emph{Convex Optimization} book (\url{https://web.stanford.edu/\~boyd/cvxbook/bv_cvxbook.pdf}):
	% \begin{itemize}
	% 	\item \textbf{Exercises:} 3.7, 3.10, 3.11 (first part), 3.21 (lots of other problems if you want more practice, but many are on the harder side)
	% \end{itemize}
	
	% \begin{enumerate}
	% 	\item Let $f_1(x), \ldots, f_n(x)$ be $\beta$-smooth convex functions and let $g(x) = \frac{1}{n}\sum_{i=1}^n f_i(x)$ be their average. Show that $g$ is $\beta$-smooth.
	% 	\item Let $f: \R\rightarrow \R$ be a $\beta$-smooth, $\alpha$-strongly  convex function. Let $g(x) = f(c\cdot x)$ for some constant $0< c < 1$.  How does $g$'s smoothness and strong convexity compare to that of $f$? How about $g$'s condition number?
	% 	%	\item We stated in class that a $\beta$-smooth convex function can be optimized to error $\epsilon$ in $O(\beta R^2/\epsilon)$ iterations of gradient descent. Suppose $f$ is not $\beta$-smooth -- we only know it is $G$-Lipschitz. 
	% 	\item Let $f(x) = x^4$. Is $f$ $G$-Lipschitz for finite $G$?  Is $f$ $\beta$-smooth for finite $B$? 
	% 	%	\item In stochastic gradient descent, we have a function $f(\bv{x}) = \sum_{i=1}^n f_i(\bv{x})$. At each step we pick an $f_i$ uniformly at random and update $\bv{x}^{(t+1)} = \bv{x}^{(t)} - \eta \nabla f_i(\bv{x}^{(t)} )$. In practice a different strategy is often used due to cache considerations: instead of picking $f_i$ at random, we cycle through functions $f_1, \ldots, f_n$, using each of their gradients in order. Specifically, we update $\bv{x}^{(t+1)} = \bv{x}^{(t)} - \eta \nabla f_j(\bv{x}^{(t)})$ where $j = t\mod (n+1)$.
	% 	%	
	% 	%	Show that this strategy can fail terribly in the worst case. In particular, describe a function $f(\bv{x}) = \sum_{i=1}^n f_i(\bv{x})$ where cyclic updates 
	% \end{enumerate}
	
	%	\item 
	%	
	%	ALWAYS\hspace{1em} SOMETIMES\hspace{1em} NEVER\vspace{4em}
	
\end{document}